% Generated draft: requires mathematical and source review.
\documentclass[11pt]{article}
\usepackage[margin=25mm]{geometry}
\usepackage{fontspec}
\setmainfont{Latin Modern Roman}
\usepackage{amsmath,amssymb,mathtools,unicode-math}
\setmathfont{Latin Modern Math}
\usepackage{enumitem,needspace,xurl,hyperref}
\setlength{\parindent}{0pt}
\setlength{\parskip}{5pt}
\setlength{\emergencystretch}{3em}
\newcommand{\N}{\mathbb{N}}
\newcommand{\Z}{\mathbb{Z}}
\newcommand{\Q}{\mathbb{Q}}
\newcommand{\R}{\mathbb{R}}
\newcommand{\C}{\mathbb{C}}
\newcommand{\F}{\mathbb{F}}
\newcommand{\A}{\mathbb{A}}
\newcommand{\PP}{\mathbb{P}}
\newcommand{\E}{\mathbb{E}}
\newcommand{\eps}{\varepsilon}
\newcommand{\dd}{\,\mathrm{d}}
\begin{document}
\section*{UnsolvedMath ID: 2308024}
\textbf{Source problem code:} catalogue entry
\subsection{Definitions and mathematical statement}
Let $(M_k)_{k\geq0}$ be positive real numbers with $M_0=1$. The source's malformed parenthesized expression is naturally read as the binomial condition
\[
M_{k+j}\geq\binom{k+j}{j}M_kM_j\quad(k,j\geq0),
\qquad \sum_{k=0}^\infty\frac{M_k}{M_{k+1}}<\infty.
\]
This binomial reading requires verification. Let $X\subset\C$ be compact. The phrase ``rational functions with poles of $X$'' appears to mean poles \emph{off} $X$, since poles on $X$ would obstruct the displayed finite derivative norm. Under that explicit conditional reading, consider rational functions $r$ holomorphic on a neighborhood of $X$ for which
\[
\|r\|_M=\sum_{k=0}^\infty\frac{\sup_{z\in X}|r^{(k)}(z)|}{M_k}<\infty,
\]
and let $A_M(X)$ be their completion in this norm, with the source's asserted Banach-algebra multiplication. Elements are the intended norm-limits of rational functions. For general compact $X$, whether the completion embeds injectively as functions on $X$, rather than carrying additional derivative data, also requires the source's conventions.

For nonempty $X$, evaluation at each $x\in X$ extends continuously to $A_M(X)$ because $|r(x)|\leq\|r\|_M$. Its maximal ideal space is the space of nonzero continuous multiplicative complex linear functionals. A. G. O'Farrell asks whether the evaluation map
\[
X\longrightarrow\mathfrak M(A_M(X)),\qquad x\longmapsto\operatorname{ev}_x
\]
is a homeomorphism, equivalently whether every character is evaluation at a unique point of $X$. The source reports yes for the unit disc and unit interval; since $X$ is compact, the disc case is naturally the closed unit disc, but that boundary convention should be checked. The supplied summary reports openness as of 2018 and no definitive resolution located.
\subsection{Short English statement}
The question asks whether a weighted derivative-norm algebra generated by rational functions has only point-evaluation characters. The sequence inequality and pole location are malformed in the extraction, and the completion's interpretation on arbitrary compact sets needs checking before a definitive formulation is fixed.
\subsection{Sources}
\begin{itemize}
\item Local UnsolvedMath registry, permanent ID 2308024, exact source code catalogue entry, ``Research Problems in Function Theory --- Problem 8.24''; status open, classification OPEN-TRIAGE.
\item Hayman, \emph{Research Problems in Function Theory} (2018), Problem 8.24, attributed to A. G. O'Farrell, \url{https://arxiv.org/abs/1809.07200}. The record identifies extraction from public arXiv source TeX, accessed July 29, 2026; the 2018 update was retained for status triage.
\item The source reports the unit-disc and unit-interval cases affirmative without a separate citation. It calls the weighted norm-limit class a Banach algebra; its exact completion/function conventions are not reproduced.
\item The supplied research summary reports openness as of the 2018 edition and no definitive resolution located, with no additional authoritative reference or precise search date.
\end{itemize}
\end{document}
